\begin{resumo}[RESUMO]
Este trabalho apresenta uma revisão do plano de crescimento e escalabilidade da startup PetCuida, plataforma digital voltada à prevenção e à ampliação do acesso a cuidados para animais de companhia. A análise documental situa o negócio em validação inicial: há protótipo funcional e evidência exploratória de usabilidade, mas ainda não há comprovação de receita, retenção, aquisição recorrente ou disposição a pagar. Com base na orientação da Entrega XIII, nas atividades anteriores e em literatura acadêmica sobre plataformas bilaterais, experimentação empreendedora e escalabilidade, propõe-se priorizar a densidade local por meio de parcerias e testar, de forma complementar, a melhoria da receita por usuário. O plano adota piloto assistido de oito semanas, indicadores com definições operacionais, critérios de decisão e expansão por etapas. A simulação financeira é tratada como análise de cenários, não como previsão: a contribuição unitária estimada de R\$ 0,55/mês e o custo do fundo de créditos indicam que o crescimento de usuários, isoladamente, não assegura sustentabilidade. A revisão também explicita riscos técnicos, operacionais, regulatórios e de proteção de dados, além de corrigir a distinção entre metas, hipóteses e resultados observados.
\vspace{\onelineskip}
\noindent\textbf{Palavras-chave}: startup; crescimento; escalabilidade; plataforma bilateral; experimentação; economia unitária.
\end{resumo}